% Umschaltung zwischen Simulation und Realbetrieb unterhalb der Herstellerschnittstelle
% (unterer Teil der früheren Gesamtarchitektur; Teilung auf Wunsch Nils 23.11.2026).
% Eigene Darstellung (Art E, Kap. 6.2.2). Keine Quelle in der Caption.
% Basis: G1@25d5c1a g1pilot/docs/02-T-Architektur.md („Sim <-> Real: Ein Schalter“, Laufzeit-Topologie),
%   g1pilot/utils/common.py:21-23 (Domain 1/lo gegen 0/Netzwerkkarte), docs/22-T-Inspire-Haende.md (Hände real über Modbus TCP).
% Stil: ba-grafik.tex. Einbinden: % Umschaltung zwischen Simulation und Realbetrieb unterhalb der Herstellerschnittstelle
% (unterer Teil der früheren Gesamtarchitektur; Teilung auf Wunsch Nils 23.11.2026).
% Eigene Darstellung (Art E, Kap. 6.2.2). Keine Quelle in der Caption.
% Basis: G1@25d5c1a g1pilot/docs/02-T-Architektur.md („Sim <-> Real: Ein Schalter“, Laufzeit-Topologie),
%   g1pilot/utils/common.py:21-23 (Domain 1/lo gegen 0/Netzwerkkarte), docs/22-T-Inspire-Haende.md (Hände real über Modbus TCP).
% Stil: ba-grafik.tex. Einbinden: % Umschaltung zwischen Simulation und Realbetrieb unterhalb der Herstellerschnittstelle
% (unterer Teil der früheren Gesamtarchitektur; Teilung auf Wunsch Nils 23.11.2026).
% Eigene Darstellung (Art E, Kap. 6.2.2). Keine Quelle in der Caption.
% Basis: G1@25d5c1a g1pilot/docs/02-T-Architektur.md („Sim <-> Real: Ein Schalter“, Laufzeit-Topologie),
%   g1pilot/utils/common.py:21-23 (Domain 1/lo gegen 0/Netzwerkkarte), docs/22-T-Inspire-Haende.md (Hände real über Modbus TCP).
% Stil: ba-grafik.tex. Einbinden: % Umschaltung zwischen Simulation und Realbetrieb unterhalb der Herstellerschnittstelle
% (unterer Teil der früheren Gesamtarchitektur; Teilung auf Wunsch Nils 23.11.2026).
% Eigene Darstellung (Art E, Kap. 6.2.2). Keine Quelle in der Caption.
% Basis: G1@25d5c1a g1pilot/docs/02-T-Architektur.md („Sim <-> Real: Ein Schalter“, Laufzeit-Topologie),
%   g1pilot/utils/common.py:21-23 (Domain 1/lo gegen 0/Netzwerkkarte), docs/22-T-Inspire-Haende.md (Hände real über Modbus TCP).
% Stil: ba-grafik.tex. Einbinden: \input{Bilder/tikz/sim-real-umschaltung} in einer figure-Umgebung. Nicht skalieren.
\begin{tikzpicture}[ba]
  % Layout für 12pt-Grundschrift der Arbeit geprüft (Vorschau immer im Arbeits-PDF ansehen).
  \tikzset{
    klein/.style={font=\footnotesize, text=baGrau},
  }
  % Herstellerschnittstelle und Modbus-Block wie in fig:architektur (Nils, 23.11.2026)
  \node[ba block, fill=baHellgrau!40, minimum height=11mm, minimum width=117mm, anchor=north] (DDS) at (7.1,4.6)
    {Herstellerschnittstelle (Unitree-DDS)\\{\footnotesize dieselben Nachrichten in Simulation und Realbetrieb}};
  \node[ba block, fill=baHellgrau!40, minimum height=11mm, minimum width=25mm, font=\footnotesize, align=center,
        anchor=north] (MB) at (14.45,4.6) {Modbus TCP\\(Realbetrieb)};
  % Umschalter
  \node[draw=baAkzent, ba normal, fill=white, rounded corners=1.5pt, text=black, align=center,
        minimum height=9mm, inner sep=3pt] (SW) at (7.6,2.6) {Umschalter\\Simulation / Realbetrieb};
  \draw[baAkzent, ba normal] (7.6,3.5) -- (SW.north);
  % Simulation
  \draw[draw=baGrau, ba duenn, fill=baHellgrau!25, rounded corners=3pt] (1.2,-1.55) rectangle (7.0,1.05);
  \node[klein, anchor=north west] at (1.3,1.0) {Container Simulation};
  \node[ba block, minimum width=0pt, text width=22mm] (SIM) at (2.75,-0.6) {MuJoCo-Simulation};
  \node[ba block, minimum width=0pt, text width=22mm] (SH) at (5.45,-0.6) {Simulierte Hände};
  % Realbetrieb
  \draw[draw=baGrau, ba duenn, fill=baHellgrau!25, rounded corners=3pt] (9.3,-1.55) rectangle (15.7,1.05);
  \node[klein, anchor=north west] at (9.4,1.0) {Realbetrieb};
  \node[ba block, minimum width=0pt, text width=30mm] (RB) at (11.4,-0.6) {Unitree G1 mit\\Onboard-Regler};
  \node[ba block, minimum width=0pt, text width=20mm] (RH) at (14.45,-0.6) {Inspire-\\Hände};
  % Verbindungen über den Umschalter
  \draw[baAkzent, ba normal, <->] (SW.west) -| (3.6,1.05);
  \draw[baAkzent, ba normal, <->] (SW.east) -| (12.4,1.05);
  % Beschriftungen innen neben den senkrechten Pfeilen, unterhalb des Umschalters (kreuzen nichts)
  \node[font=\footnotesize, anchor=north west, align=left, inner sep=1pt] at (3.75,2.05) {Domain 1,\\Loopback};
  \node[font=\footnotesize, anchor=north east, align=right, inner sep=1pt] at (12.25,2.05) {Domain 0,\\Netzwerkkarte};
  % Hände real: eigener Anschluss über Modbus TCP, an der Herstellerschnittstelle vorbei
  \draw[ba normal, <->] (MB.south) -- (RH.north);
\end{tikzpicture}
 in einer figure-Umgebung. Nicht skalieren.
\begin{tikzpicture}[ba]
  % Layout für 12pt-Grundschrift der Arbeit geprüft (Vorschau immer im Arbeits-PDF ansehen).
  \tikzset{
    klein/.style={font=\footnotesize, text=baGrau},
  }
  % Herstellerschnittstelle und Modbus-Block wie in fig:architektur (Nils, 23.11.2026)
  \node[ba block, fill=baHellgrau!40, minimum height=11mm, minimum width=117mm, anchor=north] (DDS) at (7.1,4.6)
    {Herstellerschnittstelle (Unitree-DDS)\\{\footnotesize dieselben Nachrichten in Simulation und Realbetrieb}};
  \node[ba block, fill=baHellgrau!40, minimum height=11mm, minimum width=25mm, font=\footnotesize, align=center,
        anchor=north] (MB) at (14.45,4.6) {Modbus TCP\\(Realbetrieb)};
  % Umschalter
  \node[draw=baAkzent, ba normal, fill=white, rounded corners=1.5pt, text=black, align=center,
        minimum height=9mm, inner sep=3pt] (SW) at (7.6,2.6) {Umschalter\\Simulation / Realbetrieb};
  \draw[baAkzent, ba normal] (7.6,3.5) -- (SW.north);
  % Simulation
  \draw[draw=baGrau, ba duenn, fill=baHellgrau!25, rounded corners=3pt] (1.2,-1.55) rectangle (7.0,1.05);
  \node[klein, anchor=north west] at (1.3,1.0) {Container Simulation};
  \node[ba block, minimum width=0pt, text width=22mm] (SIM) at (2.75,-0.6) {MuJoCo-Simulation};
  \node[ba block, minimum width=0pt, text width=22mm] (SH) at (5.45,-0.6) {Simulierte Hände};
  % Realbetrieb
  \draw[draw=baGrau, ba duenn, fill=baHellgrau!25, rounded corners=3pt] (9.3,-1.55) rectangle (15.7,1.05);
  \node[klein, anchor=north west] at (9.4,1.0) {Realbetrieb};
  \node[ba block, minimum width=0pt, text width=30mm] (RB) at (11.4,-0.6) {Unitree G1 mit\\Onboard-Regler};
  \node[ba block, minimum width=0pt, text width=20mm] (RH) at (14.45,-0.6) {Inspire-\\Hände};
  % Verbindungen über den Umschalter
  \draw[baAkzent, ba normal, <->] (SW.west) -| (3.6,1.05);
  \draw[baAkzent, ba normal, <->] (SW.east) -| (12.4,1.05);
  % Beschriftungen innen neben den senkrechten Pfeilen, unterhalb des Umschalters (kreuzen nichts)
  \node[font=\footnotesize, anchor=north west, align=left, inner sep=1pt] at (3.75,2.05) {Domain 1,\\Loopback};
  \node[font=\footnotesize, anchor=north east, align=right, inner sep=1pt] at (12.25,2.05) {Domain 0,\\Netzwerkkarte};
  % Hände real: eigener Anschluss über Modbus TCP, an der Herstellerschnittstelle vorbei
  \draw[ba normal, <->] (MB.south) -- (RH.north);
\end{tikzpicture}
 in einer figure-Umgebung. Nicht skalieren.
\begin{tikzpicture}[ba]
  % Layout für 12pt-Grundschrift der Arbeit geprüft (Vorschau immer im Arbeits-PDF ansehen).
  \tikzset{
    klein/.style={font=\footnotesize, text=baGrau},
  }
  % Herstellerschnittstelle und Modbus-Block wie in fig:architektur (Nils, 23.11.2026)
  \node[ba block, fill=baHellgrau!40, minimum height=11mm, minimum width=117mm, anchor=north] (DDS) at (7.1,4.6)
    {Herstellerschnittstelle (Unitree-DDS)\\{\footnotesize dieselben Nachrichten in Simulation und Realbetrieb}};
  \node[ba block, fill=baHellgrau!40, minimum height=11mm, minimum width=25mm, font=\footnotesize, align=center,
        anchor=north] (MB) at (14.45,4.6) {Modbus TCP\\(Realbetrieb)};
  % Umschalter
  \node[draw=baAkzent, ba normal, fill=white, rounded corners=1.5pt, text=black, align=center,
        minimum height=9mm, inner sep=3pt] (SW) at (7.6,2.6) {Umschalter\\Simulation / Realbetrieb};
  \draw[baAkzent, ba normal] (7.6,3.5) -- (SW.north);
  % Simulation
  \draw[draw=baGrau, ba duenn, fill=baHellgrau!25, rounded corners=3pt] (1.2,-1.55) rectangle (7.0,1.05);
  \node[klein, anchor=north west] at (1.3,1.0) {Container Simulation};
  \node[ba block, minimum width=0pt, text width=22mm] (SIM) at (2.75,-0.6) {MuJoCo-Simulation};
  \node[ba block, minimum width=0pt, text width=22mm] (SH) at (5.45,-0.6) {Simulierte Hände};
  % Realbetrieb
  \draw[draw=baGrau, ba duenn, fill=baHellgrau!25, rounded corners=3pt] (9.3,-1.55) rectangle (15.7,1.05);
  \node[klein, anchor=north west] at (9.4,1.0) {Realbetrieb};
  \node[ba block, minimum width=0pt, text width=30mm] (RB) at (11.4,-0.6) {Unitree G1 mit\\Onboard-Regler};
  \node[ba block, minimum width=0pt, text width=20mm] (RH) at (14.45,-0.6) {Inspire-\\Hände};
  % Verbindungen über den Umschalter
  \draw[baAkzent, ba normal, <->] (SW.west) -| (3.6,1.05);
  \draw[baAkzent, ba normal, <->] (SW.east) -| (12.4,1.05);
  % Beschriftungen innen neben den senkrechten Pfeilen, unterhalb des Umschalters (kreuzen nichts)
  \node[font=\footnotesize, anchor=north west, align=left, inner sep=1pt] at (3.75,2.05) {Domain 1,\\Loopback};
  \node[font=\footnotesize, anchor=north east, align=right, inner sep=1pt] at (12.25,2.05) {Domain 0,\\Netzwerkkarte};
  % Hände real: eigener Anschluss über Modbus TCP, an der Herstellerschnittstelle vorbei
  \draw[ba normal, <->] (MB.south) -- (RH.north);
\end{tikzpicture}
 in einer figure-Umgebung. Nicht skalieren.
\begin{tikzpicture}[ba]
  % Layout für 12pt-Grundschrift der Arbeit geprüft (Vorschau immer im Arbeits-PDF ansehen).
  \tikzset{
    klein/.style={font=\footnotesize, text=baGrau},
  }
  % Herstellerschnittstelle und Modbus-Block wie in fig:architektur (Nils, 23.11.2026)
  \node[ba block, fill=baHellgrau!40, minimum height=11mm, minimum width=117mm, anchor=north] (DDS) at (7.1,4.6)
    {Herstellerschnittstelle (Unitree-DDS)\\{\footnotesize dieselben Nachrichten in Simulation und Realbetrieb}};
  \node[ba block, fill=baHellgrau!40, minimum height=11mm, minimum width=25mm, font=\footnotesize, align=center,
        anchor=north] (MB) at (14.45,4.6) {Modbus TCP\\(Realbetrieb)};
  % Umschalter
  \node[draw=baAkzent, ba normal, fill=white, rounded corners=1.5pt, text=black, align=center,
        minimum height=9mm, inner sep=3pt] (SW) at (7.6,2.6) {Umschalter\\Simulation / Realbetrieb};
  \draw[baAkzent, ba normal] (7.6,3.5) -- (SW.north);
  % Simulation
  \draw[draw=baGrau, ba duenn, fill=baHellgrau!25, rounded corners=3pt] (1.2,-1.55) rectangle (7.0,1.05);
  \node[klein, anchor=north west] at (1.3,1.0) {Container Simulation};
  \node[ba block, minimum width=0pt, text width=22mm] (SIM) at (2.75,-0.6) {MuJoCo-Simulation};
  \node[ba block, minimum width=0pt, text width=22mm] (SH) at (5.45,-0.6) {Simulierte Hände};
  % Realbetrieb
  \draw[draw=baGrau, ba duenn, fill=baHellgrau!25, rounded corners=3pt] (9.3,-1.55) rectangle (15.7,1.05);
  \node[klein, anchor=north west] at (9.4,1.0) {Realbetrieb};
  \node[ba block, minimum width=0pt, text width=30mm] (RB) at (11.4,-0.6) {Unitree G1 mit\\Onboard-Regler};
  \node[ba block, minimum width=0pt, text width=20mm] (RH) at (14.45,-0.6) {Inspire-\\Hände};
  % Verbindungen über den Umschalter
  \draw[baAkzent, ba normal, <->] (SW.west) -| (3.6,1.05);
  \draw[baAkzent, ba normal, <->] (SW.east) -| (12.4,1.05);
  % Beschriftungen innen neben den senkrechten Pfeilen, unterhalb des Umschalters (kreuzen nichts)
  \node[font=\footnotesize, anchor=north west, align=left, inner sep=1pt] at (3.75,2.05) {Domain 1,\\Loopback};
  \node[font=\footnotesize, anchor=north east, align=right, inner sep=1pt] at (12.25,2.05) {Domain 0,\\Netzwerkkarte};
  % Hände real: eigener Anschluss über Modbus TCP, an der Herstellerschnittstelle vorbei
  \draw[ba normal, <->] (MB.south) -- (RH.north);
\end{tikzpicture}
